\section{From symmetry to swap operators}
\label{sec:general_formalism}

This section contains an abstract discussion of
stochastic particle systems on $\mathbb{Z}$ which
depend symmetrically on their parameters.
The main notions which we use in other parts of the
paper are \emph{parameter-symmetric stochastic particle system}
and \emph{swap operators}.

\subsection{Parameter-symmetric particle systems}
%\label{sub:symm_IPS_def}

Let $\mathrm{Conf}_{\rm fin}(\mathbb{Z})$ be the space of
particle configurations $\mathbf{x}=(\cdots <x_3<x_2<x_1)$,
$x_i\in \mathbb{Z}$,
which can be obtained from the \emph{step configuration}
$\mathsf{step}:=(\ldots,-3,-2,-1 )$
by finitely many operations of moving a particle to the right by one
into the nearby empty spot. The space
$\mathrm{Conf}_{\rm fin}(\mathbb{Z})$ is~countable.

By a multiparameter interacting particle system $\mathbf{x}(t)$
we mean a Markov process on $\mathrm{Conf}_{\rm fin}(\mathbb{Z})$ evolving in continuous
or discrete time,
such that $\mathbf{x}(0)=\mathsf{step}$.
Assume that this Markov process depends on countably many parameters
${\boldsymbol \nu}=\{\nu_i\}_{i\in \mathbb{Z}_{\ge1}}$.
The parameters $\nu_i$ in our situation are real, though without loss of generality
they may belong to an abstract space.
One should think that $\nu_i$ is attached to the particle $x_i$, but the distribution
of each $x_j(t)$ may depend on all of~the~$\nu_i$'s.
We denote the process depending on ${\boldsymbol \nu}$ by $\mathbf{x}^{{\boldsymbol \nu}}(t)$.
In this section we assume that all the parameters $\nu_i$ are pairwise distinct.

The infinite symmetric group
$S(\infty)=\bigcup_{n=1}^{\infty}S(n)$
acts on the parameters ${\boldsymbol \nu}$ by permutations,
$\sigma\colon {\boldsymbol \nu}\mapsto \sigma{\boldsymbol \nu}$.
Here $S(n)$ is the symmetric group which permutes only the
first parame\-ters~$\nu_1,\ldots,\nu_n$.
Let us denote by $S_n(\infty)\subset S(\infty)$
the subgroup which permutes $\nu_{n+1},\nu_{n+2},\ldots $
and maps each $\nu_i$, $1\le i\le n$, into itself. Note that
$S(n)\cap S_n(\infty)=S(n)\cap S_{n-1}(\infty)=\{e\}$,
and~$S(n+1)\cap S_{n-1}(\infty)=\{e,s_{n}\}$,
where $e$ is the identity permutation, and
$s_{n}=(n,n+1)$ is the transposition
$n\leftrightarrow n+1$.

By imposing a specific distributional symmetry
of $\mathbf{x}^{{\boldsymbol \nu}}(t)$
under the action of $S(\infty)$
on the parameters ${\boldsymbol \nu}$, we arrive at the following definition:

\begin{Definition}
	\label{def:symm_IPS}
	A multiparameter particle system $\mathbf{x}^{{\boldsymbol \nu}}(t)$ is called \emph{parameter-symmetric},
	if for all $n$ and $t$ we have
	the following equality of joint distributions:
\begin{gather}
\big(\ldots,x_{n+2}^{{\boldsymbol \nu}}(t),x_{n+1}^{{\boldsymbol \nu}}(t),
x_{n-1}^{{\boldsymbol \nu}}(t),\ldots, x_1^{{\boldsymbol \nu}}(t)\big)\nonumber
\\ \qquad
{}\stackrel{d}{=}\big(\ldots,x_{n+2}^{s_n{\boldsymbol \nu}}(t),x_{n+1}^{s_n{\boldsymbol \nu}}(t),
x_{n-1}^{s_n{\boldsymbol \nu}}(t),\ldots,x_1^{s_n{\boldsymbol \nu}}(t)\big).
\label{eq:symm_IPS}
\end{gather}
	That is, the transposition $s_n$ preserves the joint
	distribution of all particles except $x_n$.
\end{Definition}
Here is a straightforward corollary of this definition:
\begin{Corollary}
	In a parameter-symmetric particle system,
	for any $t$ and any $\sigma\in S(n)\cup S_{n}(\infty)$, the
	random variables
	$x_n^{{\boldsymbol \nu}}(t)$ and
	$x_n^{\sigma {\boldsymbol \nu}}(t)$
	have the same distribution.
\end{Corollary}

